% Generated draft: requires mathematical and source review.
\documentclass[11pt]{article}
\usepackage[margin=25mm]{geometry}
\usepackage{fontspec}
\setmainfont{Latin Modern Roman}
\usepackage{amsmath,amssymb,mathtools,unicode-math}
\setmathfont{Latin Modern Math}
\usepackage{enumitem,needspace,xurl,hyperref}
\setlength{\parindent}{0pt}
\setlength{\parskip}{5pt}
\setlength{\emergencystretch}{3em}
\newcommand{\N}{\mathbb{N}}
\newcommand{\Z}{\mathbb{Z}}
\newcommand{\Q}{\mathbb{Q}}
\newcommand{\R}{\mathbb{R}}
\newcommand{\C}{\mathbb{C}}
\newcommand{\F}{\mathbb{F}}
\newcommand{\A}{\mathbb{A}}
\newcommand{\PP}{\mathbb{P}}
\newcommand{\E}{\mathbb{E}}
\newcommand{\eps}{\varepsilon}
\newcommand{\dd}{\,\mathrm{d}}
\begin{document}
\section*{UnsolvedMath ID: 30000820}
\textbf{Source problem code:} catalogue entry
\subsection{Definitions and mathematical statement}
Let \(M\) be the closed aspherical three-manifold in the source, put \(\Gamma=\pi_1(M)\), \(H=H_1(M;\Z)/\mathrm{torsion}\), and \(\Lambda=\Z[H]\). The task is to construct one noncommutative Alexander invariant \(\widetilde\Delta(M)\) satisfying all the following requirements simultaneously.
\begin{enumerate}
\item Its abelianization is the classical Alexander polynomial. Since \(H\) is free abelian, \(\Lambda\) itself is commutative; the noncommutative ambient ring, group or \(K\)-theoretic object for \(\widetilde\Delta\), and the abelianization map to \(\Lambda\), are not specified. The classical polynomial's unit ambiguity and any lift normalization must be fixed, not silently replaced by exact equality of chosen representatives.
\item It satisfies the natural involution symmetry. On a group ring the standard involution sends \(g\) to \(g^{-1}\) and reverses multiplication; on \(\Lambda\) it sends \(h\) to \(h^{-1}\). The intended symmetry, including unit factors, side of multiplication and coefficients, is not stated, so no literal self-equality is inserted.
\item Its support width computes the Thurston norm. If the invariant were a nonzero finite group-ring expression \(\sum_g a_g g\), a usual directional width would be
\[
 \max_{a_g\ne0}\xi(g)-\min_{a_g\ne0}\xi(g),
\]
where \(\xi:\Gamma\to\R\) factors through \(H\). The source does not specify such a finite expression or a support/polytope construction for a more general invariant. In the oriented smooth convention, the Thurston norm on integral cohomology is the minimum of \(\sum_j\max(0,-\chi(S_j))\) over embedded surfaces representing the Poincar\'e-dual class, extended homogeneously and continuously to real cohomology. The requested equality is exact in the intended normalizations; no unspecified correction term or scalar is supplied.
\item For the intended classes \(\xi\), existence of a nonsingular closed one-form \(\omega\) with \([\omega]=\xi\) is equivalent to invertibility of \(\widetilde\Delta(M)\) in the corresponding Novikov completion. Nonsingular means \(\omega_x\ne0\) at every point, not merely that \(\xi\ne0\). The original also calls this equivalent to an extremal support coefficient condition; that condition is not printed and is retained as an additional claimed equivalence requiring definition, not discarded.
\end{enumerate}
For reference, a common positive-direction group-ring Novikov completion permits sums \(\sum_g a_gg\) with only finitely many nonzero terms satisfying \(\xi(g)\leq C\), for every real \(C\). The opposite sign convention gives a different displayed completion; the source does not choose one or specify an alternative ambient algebra. Invertibility must mean a unit in that chosen ring, with a two-sided inverse where noncommutative. It is not merely nonzero or invertible after an unmentioned field extension.

The source omits orientability/smoothness conventions, the domain of \(\xi\) (integral or real), its zero-class treatment, and degeneracies such as a zero Alexander invariant or empty support. No extra torsion or acyclicity hypothesis is guessed. In particular \(\widetilde\Delta\) is a common invariant of \(M\), not separately chosen for each \(\xi\).

The literature triage checked 2026-08-21 reports no general construction or counterexample with all the requested properties. Although it calls them three, the explicit statement includes abelianization, symmetry, width and invertibility requirements; all are retained. The open label and reported source verification remain dated metadata, not independently checked current results.
\subsection{Short English statement}
Construct a single noncommutative Alexander invariant with the prescribed classical image, involution, exact Thurston-width formula and Novikov-unit test for nonsingular closed one-forms. The original extremal-coefficient equivalence is also retained. The ambient algebra, normalizations and completion are missing, so no familiar invariant is substituted.
\subsection{Sources}
\begin{itemize}
\item Local UnsolvedMath registry, permanent ID 30000820, source code catalogue entry, ``Noncommutative Alexander Polynomials and the Thurston Norm,'' including all explicit properties and the original extremal-support-coefficient equivalence omitted from the cleaned sentence.
\item \emph{Mini-Workshop: Topology of closed one-forms and Cohomology Jumping Loci (2008)}, Oberwolfach Reports, pp.\ 2321--2360, DOI 10.4171/owr/2007/40: \url{https://doi.org/10.4171/owr/2007/40}. This is the sole supplied literature reference. The citation uses 2008, while the DOI and proposed-year metadata use 2007.
\item Literature triage checked 2026-08-21; record updated 2026-09-16. No external source has been visited or independently checked here.
\end{itemize}
\end{document}
